Beginne unten: Der untere Stab hängt waagrecht, wenn die Drehmomente bezüglich seiner Schnur gleich gross sind ($g$ kürzt sich weg):
\begin{align}
 m_2\,g\,b_1 &= m_3\,g\,b_2 \\
 m_3 &= \mKcF \\
   &= \frac{\mKb\cdot\bLi}{\bRe} \\
   &= \mKc \approx \result{\mKcgP{0}{2}}
\end{align}
Der obere Stab trägt rechts den ganzen unteren Teil, also die Gewichtskraft von $m_2+m_3$ (die Stäbe sind leicht):
\begin{align}
 m_1\,g\,a_1 &= (m_2+m_3)\,g\,x \\
 x &= \xSF \\
   &= \frac{\mKa\cdot\aOb}{\mKb+\mKc} \\
   &= \xS \approx \result{\xScmP{0}{2}}
\end{align}
